\documentclass[11pt,a4paper]{article}
\usepackage[T1]{fontenc}
\usepackage[utf8]{inputenc}
\usepackage[english]{babel}
\usepackage{amsmath,amssymb}
\usepackage[margin=2cm]{geometry}
\setlength{\parindent}{0pt}
\setlength{\emergencystretch}{3em}
\allowdisplaybreaks

\title{About \texttt{book\_exponentials\_logarithms\_beyond.tex}}
\author{}
\date{}
\begin{document}
\maketitle
This note explains what the test data in \texttt{book\_exponentials\_logarithms\_beyond.tex} is used for. That file is input for the engine in \texttt{main.c}; it is read by the test \texttt{test\_exponentials\_logarithms.c} in this folder.

\begin{enumerate}
\item \textbf{Purpose.} The file holds 6 true statements of the same kind as \texttt{book\_exponentials\_logarithms.tex} that the engine does not decide yet.

\item \textbf{What it exercises.} Honesty about the limits: the engine may answer that it cannot decide, but must never call a true statement false. When one of these becomes provable it should move to \texttt{book\_exponentials\_logarithms.tex}.

\item \textbf{Variety.} An exponent that is not of first degree ($4^{\sqrt{x+2}}$, $e^{6/(\log_2(3+x) - 2)}$), the variable in the index of a root or in base and exponent together, and three domains where the independent check with numbers cannot confirm a row of the table (a nested logarithm, or numbers too large, at a point where a radicand is exactly zero).

\item \textbf{Expected result.} The line \texttt{PASS: 0 of 6 statements \ldots are PROVED (at least 0 wanted), 0 are REFUTED (none allowed)}.

\item \textbf{How to run it.} From the top folder of the repository:
\begin{verbatim}
cc -std=c11 -O2 -Wall -Wextra -o /tmp/test_exponentials_logarithms \
    tests/exponentials_logarithms/test_exponentials_logarithms.c \
    && /tmp/test_exponentials_logarithms
\end{verbatim}
\end{enumerate}
\end{document}
